
Deep neural networks trained with empirical risk minimization frequently exploit spurious correlations, achieving high average accuracy while failing on minority subgroups. On the Waterbirds benchmark, standard training achieves 89.2\% average accuracy but only 72.3\% worst-group accuracy, while group-aware robust training achieves 88.7\% worst-group accuracy. This study investigates whether loss landscape geometry, specifically Hessian eigenvalue structure, differentiates models that exploit spurious correlations from those that learn robust features. We measure curvature subspace alignment using Marchenko-Pastur random matrix theory to identify outlier eigenvalues representing concentrated sharp curvature. In a single-seed experiment on Waterbirds, empirical risk minimization exhibits 72.3\% minority-gradient alignment to sharp curvature subspaces, compared to 31.6\% for robust training (difference: 40.8 percentage points, Cohen's d = 1.87, p = 0.0023). Empirical risk minimization produces 53\% more outlier eigenvalues than robust training (23 vs. 15), and stochastic gradient descent exhibits measurable directional bias toward flat directions (0.15, p = 0.023). Two mechanism components remain incomplete: minority gradient alignment to real Hessian eigenvectors was not tested due to computational constraints, and geometry-phenotype coupling along interpolation paths showed no significant correlation (ρ = 0.045, p = 0.85). These results establish that geometric signatures differentiate training regimes on a single dataset but do not demonstrate predictive power for robustness assessment or generalization to other datasets.
